\documentclass[a4paper,fleqn]{cas-dc}

\usepackage[numbers,sort&compress]{natbib}

%%%Author definitions
\def\tsc#1{\csdef{#1}{\textsc{\lowercase{#1}}\xspace}}
\tsc{WGM}
\tsc{QE}
\tsc{EP}
\tsc{PMS}
\tsc{BEC}
\tsc{DE}
%%% 

\begin{document}
\let\ref\Cref 		
\let\eqref\Cref 	
\let\autoref\Cref 	
\let\WriteBookmarks\relax
\def\floatpagepagefraction{1}
\def\textpagefraction{.001}
\shorttitle{R. Tirumalasetti et~al. / Engineering Science and Technology, an International Journal}

\title [mode = title]{A Smart and Efficient Online Doctor Appointment Booking System for Streamlined Healthcare Access}

\author[1]{Roopa Tirumalasetti}[type=editor, 
                        prefix=Dr.,
                        role=Supervisor]
\cormark[1]
\ead{roopa@srmap.edu.in}

\address[1]{SRM University-AP, Andhra Pradesh, India}

\author[1]{Mahesh Babu Nettem}
\author[1]{Pavan Reddy Narva}
\author[1]{Siva Karthik Valiveti}

\cortext[cor1]{Corresponding author.}

\begin{abstract} 
In the present work, an automated system of classification of X-rays of the chest detected by deep learning as a way of identifying COVID-19 and associated pulmonary disorders is introduced. The model is tested on the publicly available COVID-19 Radiography Database on Kaggle, which consists of 42,300 labelled images in four classes: COVID-19, lung opacity, normal, and viral pneumonia. To increase prediction accuracy, a test-time augmentation policy is used whereby, nine different versions of each input image, such as brightness, contrast, flipping, and sharpening transformations are created and tested, and the results are averaged. Moreover, a custom contrast enhancement algorithm based on adaptive histogram equalization is implemented to enhance the appearance of subtle radiographic appearances like ground-glass opacities that are typically related to COVID-19. An optimization bias correction method in the form of a grid-search is also presented to handle the bias of the classes and systematic misclassification, leading to an optimal prediction correction vector. The overall classification accuracy of the proposed approach is 91.75 percent, with the class-wise accuracies of 92 percent of COVID-19, 93 percent of lung opacity, 92 percent of normal cases, and 90 percent of cases of viral pneumonia. These findings show that the incorporation of preprocessing improvements and inference-level optimization strategies can strongly improve model robustness and diagnostic results on large-scale radiographic data.
\end{abstract}

\begin{keywords}
COVID-19 detection \sep 
chest X-ray classification \sep 
test-time augmentation \sep 
contrast enhancement \sep 
bias correction \sep 
deep learning
\end{keywords}

\maketitle

\section{INTRODUCTION}
The COVID-19 pandemic immediately spread around the world, which resulted in a critical requirement of effective, scalable, and efficient diagnostic systems that can help healthcare professionals to detect and monitor the disease early. However, reverse transcription polymerase chain reaction (RT-PCR) testing has become the method of choice in diagnoses, but has a number of practical issues, such as long processing time, lack of availability in resource-limited environments, and sensitivity problems in early infection. Consequently, medical imaging, especially chest radiography, has become a significant complementary diagnostic method to detect pulmonary abnormalities related to COVID-19 and other respiratory diseases. The use of chest X-ray imaging is relatively cheap, easily available and can offer quick diagnostic information, which is suitable when screening a large population.

Recent developments in artificial intelligence and deep learning have enhanced considerably the accuracy of automated systems in analyzing medical images. Convolutional neural networks have been widely used in the image classification process and have been highly successful in identifying different thoracic diseases in chest X-rays. Nonetheless, even with these improvements, there are still some issues with creating strong and generalizable models of multi-class classification of COVID-19, lung opacity, normal cases, and viral pneumonia. These problems are variability of image quality, differences in acquisition conditions, subtle visual differences between classes and inherent class imbalance in datasets.

Among the major challenges in the process of X-ray analysis of the chest, there are small and diffuse patterns (ground-glass opacities) that are typical of COVID-19 but can hardly be identified because of the low contrast and overlapping anatomy. Conventional preprocessing methods might not be effective enough to boost these weak features, resulting in low model sensitivity. Moreover, even the models that are trained using large datasets can be biased against a specific set of classes, which will lead to a systematic under- or over-prediction of particular conditions. These biases may have a substantial effect on clinical reliability and have to be overcome to provide equal performance on all categories.

In order to address these shortcomings, this paper proposes a more advanced deep learning-based classification system which incorporates several preprocessing and inference-level optimization strategies. The model is trained and tested on the COVID-19 Radiography Database on Kaggle, including 42,300 chest X-ray images and four classes, i.e., COVID-19, lung opacity, normal, and viral pneumonia. The large scale data used will guarantee the model acquires a wide range of patterns and enhances its capacity to generalize to a variety of imaging conditions.

One of the contributions of this research is that it included the idea of test-time augmentation in the inference stage. Prediction does not utilize a single input image but rather the image is transformed into a variety of differences, such as variations in brightness, contrast, orientation, and sharpness. The model works out each variation separately and the resultant prediction is obtained by averaging the outputs. This method minimizes the influence of noise and variations in input images, resulting in more consistent and accurate forecasts. Test-time augmentation is an effective method in improving the model and it does not need more training data.

As another enhancement method to be added to the preprocessing pipeline, a contrast enhancement algorithm, which is based on contrast limited adaptive histogram equalization, is applied. The approach enhances the visibility of localized structures in the chest X-ray images, which is easier to get the model to identify the key features of various pulmonary conditions. The preprocessing step helps to increase the feature extraction and classification performance by increasing subtle patterns, like ground-glass opacities.

One more significant feature of this work is the use of grid-search-based bias correction mechanism. The early assessments showed that the model tended to underestimate COVID-19 cases but overestimated other types. To solve this problem, various combinations of adjustment parameters were examined systematically, which generates an optimized bias vector that adjusts the output probabilities of the model and then the final classification is performed. The method is an effective way to overcome the imbalance of classes and enhance the equality of forecasts of all types.

The system has a test set classification accuracy of 91.75 percent and per-class accuracy of 92 per cent with COVID-19, 93 per cent with lung opacities, 92 per cent with normal cases, and 90 per cent with viral pneumonia. These findings indicate that the combination of sophisticated preprocessing, augmentation, and bias correction methods has a significant ability to improve the results of deep learning models in medical image analysis.

\section{PURPOSE OF THE PAPER}
This paper aims to create and assess an automated system to classify a chest X-ray to diagnose COVID-19 and other pulmonary diseases based on deep learning methods. The research will seek to provide solutions to some of the limitations that have been noted in the current methodologies such as the sensitivity to image quality changes, the inability to detect subtle radiographic features, and lack of balance in the model forecasts of various disease types. Based on the COVID-19 Radiography Database of Kaggle, comprising of 42,300 curated images with four labels, the suggested system concentrates on advancing the performance of classification with better preprocessing and inference methods. In particular, the paper explores the effect of test-time augmentation in prediction in which multiple transformed representations of every input image are considered to achieve a more stable output. Moreover, a contrast improvement algorithm based on an adaptive histogram equalization is added to enhance the brightness of small objects like ground-glass opacities that are often related to COVID-19. Another bias correction method, a grid-search based method, is also presented to control systematic misclassification and provide equal performance at all classes. The general goal is to develop a powerful and robust model, which is capable of helping medical image analysis by enhancing accuracy and consistency in the multi-class classification of chest X-ray, and to be applicable to a large scale real-world dataset.

\section{Literature Survey}
Deep learning in medical imaging has made a lot of progress in automated analysis of the X-ray of the chest, especially in detecting COVID-19 and other pulmonary diseases. The initial research was centered on binary classification, which classifies infected and non-infected cases, but in the recent past, multiclass classification has been gaining popularity as it is more clinically applicable. Pal et al. \cite{ref1} performed a comparative study of binary and multiclass classification frameworks between machine learning and deep learning models and indicated that multiclass frameworks offer more detailed diagnostic information although more complex. Their results emphasize the need to have a strong feature extraction and model generalization when addressing many types of diseases like COVID-19, pneumonia, and normal conditions.

Transfer learning has become a popular method to overcome the problem of limited labeled medical data. Lin and Fang \cite{ref2} investigated the idea of utilizing transfer learning models to predict ICU admission in patients with COVID-19 with regards to chest X-ray pictures. Their study revealed that trained convolutional neural networks can effectively learn the relevant radiographic features, and thus, can be effectively used to make accurate classifications when using smaller datasets. Nonetheless, the transfer learning models are still subject to domain mismatch and bias, especially when they are used on different data sets that are taken in dissimilar imaging conditions.

In a further effort to improve the classification performance, the researchers have investigated hybrid methods between the traditional image processing and deep learning. Gtifa et al. \cite{ref3} suggested a model that combines nature-inspired multi-level thresholding and convolutional neural networks to enhance the accuracy of segmentation and classification. Their method focuses on the significance of preprocessing in accentuating significant areas of interest in the chest X-ray images. Likewise, Cheluvaraju and Shivasubramanya \cite{ref5} proposed attention-based learning with fuzzy-enhanced data diversity, which optimises the model payoff on the content features, and generalisation in different image states.

Regardless of such developments, model bias and class imbalance continue to be a major issue in the medical image classification. Chen et al. \cite{ref4} have resolved this difficulty by integrating Bayesian deep learning and statistical methods to curb bias in detecting COVID-19. The results of their study showed that uncertainty estimation and minimization of systematic errors in predictions can be enhanced by probabilistic modeling. In multiclass classification, bias mitigation is especially crucial, as models can be biased towards some classes compared to others because of either imbalance in the data or similarity in the features.

The other significant field of study is to enhance model robustness in the process of inference. The use of test-time augmentation has attracted attention as a useful method to improve the stability of prediction, without changing the training process. Sherkatghanad et al. \cite{ref6} proposed the uncertainty-aware model, BayTTA, which uses test-time augmentation, and Bayesian model averaging to enhance the reliability of classification. They find that creating many augmented versions of input images and averaging predictions can greatly decrease variance and enhance accuracy. Other areas that have been considered through similar methods include medical image segmentation where Ma et al. \cite{ref7} showed that generative augmentation methods are effective at inference time.

The use of test-time augmentation in distributed and real-world applications has also been explored recently. Melo et al. \cite{ref8} tested its application in federated learning scenarios in brain tumor classification and demonstrated that test-time augmentation can improve model performance when working with decentralized data. To enhance the interpretability and diagnostic confidence of skin lesion classification, Cino et al. \cite{ref9} used test-time augmentation and explained artificial intelligence methods. Chen et al. \cite{ref10} also took this idea a step further by suggesting test-time adaptation techniques, which can optimize model predictions without the need to update parameters, and that inference-level optimization strategies are flexible.

Besides the augmentation methods, image enhancement methods are also very important in enhancing the quality of medical images and enabling accurate extraction of features. Ahmad et al. \cite{ref11} assessed the influence of contrast-limited adaptive histogram equalization and other image enhancement methods on the segmentation of X-ray images and revealed that the models were greatly enhanced in their performance. Most notably, CLAHE has been known to boost local contrast with the ability to maintain valuable structure information. Mamatov and others \cite{ref12} developed an extensive discussion of different contrast enhancement algorithms, and their relevance in medical imaging studies.

There have also been noise reduction and hybrid enhancement strategies to further enhance the image quality. Ismael et al. \cite{ref13} came up with a technique that involves noise reduction and contrast enhancement to give clearer and more informative medical images. They find that pipelines which are preprocessed by combining a variety of enhancing methods can greatly enhance the results of classification. Pertiwi et al. \cite{ref14} applied CLAHE to lung X-ray images of COVID-19 patients in particular and revealed its ability to improve visibility of critical features. On the same note, Risma and Utami \cite{ref15} emphasized the significance of contrast enhancement in enhancing the quality of deep learning-based diagnosis of tuberculosis, which supports the significance of preprocessing in medical image analysis.

Although there has been a great improvement in creating deep learning models to classify chest X-rays, various studies have indicated that there are still a number of limitations that warrant the creation of better methodologies. Among the major issues that can be observed throughout the literature is the inability to differentiate the visually similar pulmonary conditions, especially COVID-19, lung opacity, and viral pneumonia. Pal et al. \cite{ref1} pointed out that multiclass classification causes greater inter-class similarity, which may decrease the accuracy of the models in comparison with binary classification. The problem is also complicated by overlapping radiographic appearances where the patterns of opacities and consolidations can be seen in a variety of disease types.

Transfer learning methods, as effective as they are, have some limitations as well. Lin and Fang \cite{ref2} performed well with pretrained models, but their experiment showed that the model predictions are subject to biases that are transferred to the source datasets. This may cause discrepancies upon the use of the models on new datasets whose distribution is different. Consequently, there is an increasing desire to have the means that can adapt to variability of data and enhance generalization without necessarily depending on training-stage adjustments.

To overcome these limitations hybrid preprocessing and feature enhancement techniques have been researched. Gtifa et al. \cite{ref3} demonstrated that thresholding techniques can be combined with convolutional neural networks to enhance the accuracy of the segmentation, but these methods frequently need parameter fine-tuning and may not be well-generalized to a large variety of data sets. On the same note, attention based models introduced by Cheluvaraju and Shivasubramanya \cite{ref5} are able to enhance localization of features at the cost of increasing additional computational complexity, which could be restricting to real-time applications.

Biases in model forecasting are a pressing issue, and especially in healthcare where misclassification may be very costly. Chen et al. \cite{ref4} showed that Bayesian deep learning models are able to estimate uncertainties more effectively and minimize prediction bias. Nonetheless, they tend to be associated with higher computational costs, and they can be associated with complicated training processes. This shows that easier but effective bias correction mechanisms that can be used in inference are necessary.

Test-time augmentation has proved as a potential solution to enhance model robustness without any changes in training process. Sherkatghanad et al. \cite{ref6} proposed an uncertainty-contained augmentation framework, which averages multiple predictions with the help of the Bayesian framework, leading to better classification rates. Their results indicate that inference-level methodologies can be successfully used to overcome constraints of training data. Ma et al. \cite{ref7} also investigated generative augmentation methods, and they showed that synthetic variations of input images can increase segmentation accuracy. These methods show that by using multiple representations of the same input the stability of the model can be enhanced.

Test-time augmentation is not only applicable in centralized learning systems. Melo et al. \cite{ref8} explored its application in a federated learning setting, where data is spread among several nodes. Their research revealed that the test-time augmentation is capable of enhancing model performance despite the fact that training data is not centrally available. Equally, Cino et al. \cite{ref9} used augmentation alongside clarifiable artificial intelligence in the assessment of skin lesions to augment accuracy and interpretability. Chen et al. \cite{ref10} also suggested test-time adaptation procedures that adapt predictions in real-time without modifying the model parameters, which also shows the flexibility of inference-based optimization methods.

The quality of images has been a key element of enhancing the model performance in medical imaging. Ahmad et al. \cite{ref11} emphasised the usefulness of contrast enhancement methods including CLAHE in enhancing the segmentation and classification results. Their experiment showed that the visibility of significant features can be greatly improved by increasing the local contrast, especially in low quality X-ray images. Mamatov et al. \cite{ref12} have given a thorough discussion of contrast enhancement algorithms, and their use in dealing with imaging condition variability.

Besides enhancing contrast, noise reduction methods have also been investigated to enhance the clarity of images. Ismael et al. \cite{ref13} suggested a hybrid method, using noise reducing and contrast enhancing techniques, which led to the better image quality and classification accuracy. Pertiwi et al. \cite{ref14} applied CLAHE directly to the COVID-19 chest X-ray images, showing that it was able to increase the subtle features related to the disease. Significance of preprocessing was also affirmed by Risma and Utami \cite{ref15} who demonstrated that deep learning performance is enhanced with contrast enhancement in the diagnosis of tuberculosis.

In spite of these developments, a gap in the literature concerning the combination of preprocessing, augmentation, and bias correction into one single framework has been identified. The vast majority of the current literature concentrates on individual elements, e.g., enhancement or augmentation, and does not dwell upon their interaction to increase model performance. Moreover, most solutions are based on training-phase adjustments, which in large-scale or real-time use might not be feasible.

The present study will overcome these limitations by integrating several methods into a clearly defined framework. The proposed methodology will enhance the accuracy and robustness of the classification by adding contrast enhancement, test-time augmentation, and bias correction without adding too much complexity to the computation. The method is especially applicable to large datasets like the COVID-19 Radiography Database, where image quality and the distribution of classes may vary considerably and affect the performance of the model.

\section{DATASET}
The dataset utilized in this paper is the COVID-19 Radiography Database, which was downloaded on Kaggle, the open-source database of chest X-ray images that were created to facilitate research in automated disease detection. The dataset includes 42,300 images, which are classified into four different classes: COVID-19, lung opacity, normal and viral pneumonia. All images have a label based on the medical condition and thus they can undergo supervised learning of multi-class classification problems. The presence of a great number of labeled samples in various categories renders this dataset appropriate to train and evaluate the deep learning models in pulmonary disease detection.

The images in the dataset are also taken in different sources, and hence, they have different resolutions, contrast and imaging conditions. This variability is a representation of the real-life conditions where the chest X-rays are obtained with the help of different devices and settings. Although such diversity enhances the generalization ability of trained models, it also presents the problem of uneven image quality and representation of features. Specifically, the presence of subtle patterns related to COVID-19, including ground-glass opacities, can not be easily seen in all images as a result of exposure and noise variations. Thus, the input data needs to be preprocessed with proper techniques to standardize data and improve significant features before model inference.

The data set consists of four clinically significant classes which tend to be hard to differentiate because of similar radiographic features. COVID-19 and viral pneumonia share the patterns of opacities, infiltrates, and lung opacity and normal cases may have dissimilar intensities and be classified as misclassified. Such multi-class structure makes classification task more complex and demands the model to understand fine-grained distinctions between classes. Also, the number of images in each of the classes might not be evenly distributed, generating bias in model predictions unless it is managed.

To evaluate the model proposed experimentally, the dataset serves as the main source of testing the model performance. The model works with every picture by running an augmented inference pipeline that considers a number of augmented variations of the input and can make more robust predictions. The size of the dataset is large, which guarantees that the results of the evaluation are statistically sound and reflect the actual performance. The study will use a detailed and versatile dataset to show how the methods proposed will enhance the accuracy and robustness of classification in the analysis of chest X-rays.                  

\section{PROPOSED METHODOLOGY}

\subsection{Overview Of The Proposed System}
The proposed system is designed to perform automated classification of chest X-ray images into four categories: COVID-19, lung opacity, normal, and viral pneumonia. The methodology focuses on improving prediction robustness and classification accuracy through enhanced preprocessing and inference-level optimization rather than complex model restructuring. The system takes a chest X-ray image as input and processes it through a sequence of steps including image enhancement, test-time augmentation, model inference, and bias correction. The final output is generated by aggregating predictions from multiple transformed versions of the same image. This approach ensures that the system remains stable under variations in image quality and acquisition conditions while maintaining computational efficiency.

\subsection{Input Image Processing}
The input image is first standardized to ensure compatibility with the trained deep learning model. This includes resizing the image to the required input dimensions and normalizing pixel values to a consistent scale. Since the dataset contains images from different sources with varying resolutions and lighting conditions, this step is essential for reducing variability. The preprocessing stage prepares the image for further enhancement and augmentation without altering the underlying medical features. Proper normalization also ensures that the model interprets pixel intensities consistently across all samples, which is critical for accurate feature extraction.

\subsection{Contrast Enhancement Using Clahe}
To improve the visibility of subtle radiographic patterns, a contrast enhancement technique inspired by contrast limited adaptive histogram equalization is applied. This method enhances local contrast within the image, allowing faint features such as ground-glass opacities to become more distinguishable. Unlike global histogram equalization, this approach operates on localized regions, preventing over-amplification of noise while preserving structural details. The enhanced image provides clearer input to the model, enabling it to capture important diagnostic features more effectively. This step plays a crucial role in improving classification performance, especially for COVID-19 cases where visual indicators are often subtle.

\subsection{Test Time Augmentation Strategy}
Instead of relying on a single image for prediction, the proposed system employs a test-time augmentation strategy to generate multiple variations of the input image. These variations include transformations such as horizontal flipping, brightness adjustment, contrast modification, and sharpening. A total of nine versions of the image are created, each representing a different perspective of the same input. The model processes each variation independently, producing a set of predictions. The final output is obtained by averaging these predictions, which reduces the impact of noise and inconsistencies in the original image. This approach enhances the robustness of the system without requiring additional training data or changes to the model architecture.

\subsection{Model Prediction And Aggregation}
The deep learning model used in this system performs classification on each augmented image and outputs probability scores for all four classes. These probabilities represent the model’s confidence in each category. After processing all augmented inputs, the probabilities are aggregated by computing their average across all variations. This aggregation step ensures that the final prediction is not influenced by a single image condition but reflects a consensus across multiple representations. As a result, the system produces more stable and reliable predictions, particularly in cases where individual transformations may introduce minor variations in output.

\subsection{Bias Correction Using Grid Search Optimization}
During initial evaluation, the model exhibited a tendency to under-predict COVID-19 cases and over-predict other classes. To address this issue, a bias correction mechanism is introduced using grid search optimization. Multiple combinations of adjustment parameters are tested systematically to identify an optimal bias vector. The selected bias vector is applied to the aggregated probability scores before final classification. This adjustment increases the confidence for under-represented classes and reduces overestimation for others, resulting in a more balanced prediction distribution. The bias correction step improves overall accuracy and ensures fairness across all classes without modifying the trained model.

\subsection{Final Classification Output}
After applying bias correction, the system selects the class with the highest adjusted probability as the final prediction. The output is presented as one of the four predefined categories along with its associated confidence score. This final stage ensures that the decision is based on a combination of enhanced input features, multiple augmented predictions, and corrected probability distributions. The overall methodology integrates preprocessing, augmentation, and optimization techniques into a unified pipeline, enabling accurate and robust classification of chest X-ray images in a multi-class setting.

\section{RESULTS AND DISCUSSION}

\subsection{Experimental Setup}
The proposed system was evaluated using the COVID-19 Radiography Database consisting of 42,300 chest X-ray images categorized into COVID-19, lung opacity, normal, and viral pneumonia classes. The evaluation was performed using the enhanced inference pipeline that integrates contrast enhancement, test-time augmentation, and bias correction. Each input image was processed through nine augmented variations, and predictions were aggregated to obtain the final result. The model outputs probability scores for each class, which are further refined using a bias correction vector before final classification. This setup ensures that the evaluation reflects the actual performance of the deployed system under realistic conditions.

\subsection{Overall Classification Performance}
The system achieved an overall classification accuracy of 91.75 percent on the dataset. This result indicates that the proposed methodology is capable of effectively distinguishing between multiple pulmonary conditions using chest X-ray images. The use of a large dataset contributes to the reliability of the evaluation, as it includes diverse imaging conditions and patient variations. The achieved accuracy demonstrates that combining preprocessing and inference-level optimization techniques can significantly enhance model performance without modifying the underlying architecture.

\subsection{Class Wise Performance Analysis}
The model achieved class-wise accuracies of 92 percent for COVID-19, 93 percent for lung opacity, 92 percent for normal cases, and 90 percent for viral pneumonia. These results indicate that the system performs consistently across all categories, with slightly lower performance observed for viral pneumonia. The similarity in radiographic features between viral pneumonia and other classes contributes to this variation. However, the overall distribution of accuracy values suggests that the model maintains balanced performance and does not excessively favor any single class. This balance is essential for practical medical applications where equal importance must be given to all diagnostic categories.

\subsection{Impact Of Test Time Augmentation}
The implementation of test-time augmentation plays a significant role in improving prediction stability and robustness. By generating multiple variations of each input image, the system reduces sensitivity to noise, lighting variations, and positional differences. The aggregation of predictions from nine augmented images ensures that the final output is not dependent on a single representation. This approach leads to more consistent classification results and reduces the likelihood of incorrect predictions caused by minor variations in input data. The results indicate that test-time augmentation effectively enhances the reliability of the model without requiring additional training.

\subsection{Effect Of Contrast Enhancement}
The use of contrast enhancement based on adaptive histogram equalization improves the visibility of important features within chest X-ray images. This is particularly important for detecting COVID-19, where subtle patterns such as ground-glass opacities may not be easily visible in raw images. The enhanced images provide clearer input to the model, enabling better feature extraction and classification. The observed improvement in COVID-19 detection accuracy suggests that preprocessing plays a critical role in medical image analysis. The enhancement technique ensures that important diagnostic features are preserved while minimizing the impact of noise.

\subsection{Bias Correction And Model Behavior}
The bias correction mechanism introduced in the system addresses the issue of uneven prediction distribution across classes. Initial observations indicated that the model tended to under-predict COVID-19 cases and over-predict other categories. By applying an optimized bias vector to the aggregated probability scores, the system achieves a more balanced classification outcome. This adjustment improves overall accuracy and ensures fair representation of all classes. The bias correction step is particularly important in multiclass classification scenarios, where class imbalance and feature similarity can lead to systematic errors.

\sectionCONCLUSION
The paper introduces an automated system of classifying chest X-rays to identify COVID-19 and associated pulmonary diseases based on a combination of deep learning and inference-level optimization approaches. This model was tested on the COVID-19 Radiography Database, which has 42,300 images in four classes. Overall, the system has an accuracy of 91.75 percent, which shows that it can effectively carry out reliable multi-class classification. The findings show that the model can successfully differentiate the cases of COVID-19, lung opacities, normal, and viral pneumonia against the background of the natural similarity of radiographic characteristics. A large and varied dataset is used to make sure that the assessment is based on realistic conditions, which will justify the validity of the acquired results.

The proposed methodology aims at enhancing the performance of the models via preprocessing and inference strategies instead of altering the underlying architecture. The combination of contrast enhancement enhances the appearance of fine details in chest X-ray images, which is especially necessary to recognize COVID-19 patterns. The test-time augmentation is used to enable the system to create various versions of every input image, which leads to more consistent and stable predictions. Moreover, the bias correction mechanism is used to deal with class imbalance and systematic misclassification, whereby there is balanced performance among all the categories. The combination of these methods helps to provide the strength and stability of the system as a whole.

Although good performance was achieved, there are still some limitations towards differentiating classes with overlapping features, especially the cases of viral pneumonia and the normal cases. Such difficulties help to identify the necessity of additional improvement of feature extraction and classification techniques. Further enhancements can involve the incorporation of more sophisticated deep learning architectures or other data-driven models to improve on the discrimination of similar conditions. In general, the research study indicates that a combination of preprocessing optimization and inference-level optimization is a good strategy on how to boost the performance of the chest X-ray classification system and as such, the system is appropriate in assisting medical image analysis in real-world scenarios.

\bibliographystyle{elsarticle-num} 				
\bibliography{cas-refs}

\end{document}
